% tikz-tensors-grid.tex -- the two-dimensional grid: one tensor repeated over a
% square lattice, a block of the canvas like a stack.

% \tngrid[legs, bonds=<style>/<label>]{<name>}{<columns>}{<rows>}{<style>/<label>}
% A two-dimensional network of one tensor repeated: <columns> by <rows> of it,
% each joined to its neighbours along both directions of the lattice, drawn
% turned by 45 degrees so that the lattice reads as a plane seen from above --
% a PEPS, a two-dimensional partition function. The tensor at column i, row j
% is the node <name>-i-j. `legs' gives each one an index out of the plane,
% drawn straight down and ending at <name>-i-j-down; `bonds=' puts a tensor of
% its own on every bond, half way along -- the weights of a simple update --
% turned with the lattice so that the bond meets it as it would on a chain.
% The grid is a block of the canvas, placed as a stack is -- after what is on
% the row and centred on it -- so it takes part in an equation like any stack.
\pgfkeys{/tn/grid/.is family, /tn/grid/legs/.code={\def\tn@glegs{1}},
         /tn/grid/bonds/.initial=}
\def\tn@gxy#1#2{%
  \tn@len\tn@gx{\tn@gxz+((#1)-(#2)+\tn@grows-1)*\tn@gh}%
  \tn@len\tn@gy{\the\tn@line@y+((\tn@gcols+\tn@grows)/2-(#1)-(#2)+1)*\tn@gh}}
\newcommand{\tngrid}[5][]{%
  \def\tn@glegs{0}%
  \pgfkeys{/tn/grid/.cd, bonds=, #1}%
  \pgfkeysgetvalue{/tn/grid/bonds}\tn@gbond
  \def\tn@gcols{#3}\def\tn@grows{#4}%
  \tn@len\tn@gh{\tn@gpitch*0.7071068}%
  \tn@block{2\dimexpr\tn@gap\relax}{2*\tn@gap+(#3+#4-2)/2*\tn@gh+\tn@slot/2+\tn@stub}%
  \tn@len\tn@gxz{\the\tn@cursor+\tn@slot/2}%
  \expandafter\tn@entry\tn@gbond/////\tn@end
  \let\tn@bsty\tn@sty \let\tn@blab\tn@lab
  \tn@entry#5/////\tn@end
  \foreach \tn@i in {1,...,#3}{\foreach \tn@j in {1,...,#4}{%
    \tn@gxy{\tn@i}{\tn@j}%
    \expandafter\tn@place\expandafter{\tn@sty}{#2-\tn@i-\tn@j}{\tn@gx}{\tn@gy}{\tn@lab}%
    \ifnum\tn@glegs=1
      \tn@len\tn@yb{\tn@gy-\tn@slot/2-\tn@stub}%
      \tn@line{tn edge}{\tn@gx}{\tn@gy}{\tn@gx}{\tn@yb}%
      \coordinate (#2-\tn@i-\tn@j-down) at (\tn@gx,\tn@yb);
    \fi}}%
  \foreach \tn@i in {1,...,#3}{\foreach \tn@j in {1,...,#4}{%
    \tn@gxy{\tn@i}{\tn@j}\let\tn@xa\tn@gx \let\tn@ya\tn@gy
    \ifnum\tn@i<#3
      \tn@gxy{\tn@i+1}{\tn@j}\tn@gbondto{#2}{\tn@i}{\tn@j}{a}%
    \fi
    \ifnum\tn@j<#4
      \tn@gxy{\tn@i}{\tn@j+1}\tn@gbondto{#2}{\tn@i}{\tn@j}{b}%
    \fi}}%
  \tn@gxy{#3}{1}\tn@reach{\tn@gx+\tn@slot/2}%
}
% the bond from (\tn@xa,\tn@ya) to (\tn@gx,\tn@gy), with its tensor if any
\def\tn@gbondto#1#2#3#4{%
  \tn@line{tn edge}{\tn@xa}{\tn@ya}{\tn@gx}{\tn@gy}%
  \ifx\tn@gbond\@empty\else
    \tn@len\tn@mx{(\tn@xa+\tn@gx)/2}\tn@len\tn@my{(\tn@ya+\tn@gy)/2}%
    % turned with the lattice, so that the bond meets it where it would on a
    % chain -- at a vertex of a diamond, not the middle of a side -- and its
    % label, a node of its own, upright
    \edef\tn@opt{\tn@bsty, rotate=45}%
    \expandafter\tn@place\expandafter{\tn@opt}{#1-#2-#3-#4}{\tn@mx}{\tn@my}{}%
    \edef\tn@opt{\tn@bsty, draw=none, fill=none}%
    \expandafter\tn@place\expandafter{\tn@opt}{#1-#2-#3-#4-label}{\tn@mx}{\tn@my}{\tn@blab}%
  \fi}
